\documentclass[12pt,a4paper]{article}

% Compilador: pdfLaTeX (arquivo independente)
\usepackage[utf8]{inputenc}
\usepackage[T1]{fontenc}
\usepackage[brazil]{babel}
\usepackage{lmodern}
\usepackage{geometry}
\usepackage{setspace}
\usepackage{amsmath,amssymb}
\usepackage{booktabs,array}
\usepackage{graphicx}
\usepackage{float}
\usepackage{tikz}
\usepackage{pgfplots}
\usepackage{hyperref}
\usepackage{fancyhdr}
\usepackage{needspace}
\usepackage{listings}
\usepackage{enumitem}

\pgfplotsset{compat=1.18}
\geometry{a4paper,left=3cm,right=2cm,top=2.5cm,bottom=2.5cm}
\onehalfspacing
\hypersetup{
    colorlinks=true,
    linkcolor=black,
    urlcolor=black,
    pdftitle={Identificação de Pontes e Caminhos Eulerianos},
    pdfauthor={Rafael Xavier Oliveira; Thales Duque Câmara; Gabriel Teotônio de Castro Coelho Costa}
}

\setlength{\parindent}{1.25cm}
\setlength{\parskip}{0.2em}
\pagestyle{fancy}
\setlength{\headheight}{15pt}
\fancyhf{}
\fancyhead[R]{\small Pontes e caminhos eulerianos}
\fancyfoot[C]{\thepage}
\renewcommand{\headrulewidth}{0.4pt}
\renewcommand{\footrulewidth}{0pt}
\lstset{
    basicstyle=\ttfamily\footnotesize,
    columns=fullflexible,
    breaklines=true,
    showstringspaces=false,
    numbers=none,
    frame=none
}
\setlist[itemize]{itemsep=3pt,topsep=5pt}
\newcommand{\ordem}[1]{\mathcal{O}\!\left(#1\right)}
\newcommand{\naive}{na\"ive}

\begin{document}

% =========================================================
% CAPA
% =========================================================
\begin{titlepage}
  \begin{center}
    {\large Pontifícia Universidade Católica de Minas Gerais}\\[1cm]
    {\large Ciência da Computação}\\[3.8cm]

    {\large Trabalho Prático n.º 01}\\[0.7cm]

    {\bfseries\LARGE IDENTIFICAÇÃO DE PONTES E\\[0.3cm]
    CAMINHOS EULERIANOS}\\[0.9cm]

    {\large Comparação entre os métodos naïve e de Tarjan\\
    aplicados ao algoritmo de Fleury}\\[3.4cm]

    {\large Rafael Xavier Oliveira}\\[0.3cm]
    {\large Thales Duque Câmara}\\[0.3cm]
    {\large Gabriel Teotônio de Castro Coelho Costa}\\[0.7cm]
    {\large Professor: Zenilton Kleber Gonçalves do Patrocínio Júnior}\\[2.2cm]

    {\large Belo Horizonte --- 2026}
  \end{center}
\end{titlepage}

\section{Introdução}
A modelagem por grafos permite estudar redes de comunicação, transporte e
interdependência entre elementos. Para um grafo simples não direcionado
$G=(V,E)$, com $n=|V|$ vértices e $m=|E|$ arestas, uma \emph{ponte} é uma
aresta cuja retirada aumenta o número de componentes conexas.
A identificação de pontes é especialmente relevante para o método de Fleury,
que percorre cada aresta de um grafo exatamente uma vez, quando isso é possível.

O trabalho tem três objetivos: implementar um detector de pontes por força
bruta, implementar o algoritmo linear de Tarjan e empregar os dois métodos
no algoritmo de Fleury. Para avaliar a eficiência, são utilizados grafos de
três classes eulerianas, em quatro ordens de grandeza de tamanho.
Além dos requisitos obrigatórios, o projeto oferece testes independentes,
geração reprodutível de instâncias e variantes experimentais que ajudam a
interpretar os resultados.

\section{Fundamentação teórica}
\subsection{Conectividade e pontes}
Uma aresta $e=\{u,v\}$ é ponte se, e somente se, não existe outro caminho
entre $u$ e $v$ depois da remoção de $e$. Em grafos não direcionados, isso
também equivale a dizer que $e$ não participa de nenhum ciclo.
A definição vale inclusive para grafos inicialmente desconexos: nesse caso,
a ponte é a aresta cuja remoção aumenta o número de componentes conexas.

\subsection{Critério de existência de um caminho euleriano}
Uma trilha euleriana utiliza cada aresta exatamente uma vez. Se o vértice
final coincide com o inicial, a trilha é um circuito euleriano.
Desconsiderando vértices isolados, um grafo não direcionado possui:
\begin{itemize}
  \item \textbf{Circuito euleriano:} quando é conexo e todos os vértices têm grau par;
  \item \textbf{Caminho euleriano aberto:} quando é conexo e exatamente dois vértices têm grau ímpar;
  \item \textbf{Nenhum caminho euleriano:} quando a condição de conectividade falha ou quando existem mais de dois vértices ímpares.
\end{itemize}
Se o grafo tem dois vértices ímpares, o caminho deve iniciar em um deles
e terminar no outro. A contagem dos graus e uma busca de conectividade
permitem decidir a existência de um percurso em $\ordem{n+m}$.

\subsection{Identificação de pontes pelo método naïve}
O método \naive{} considera cada aresta $e=\{u,v\}$ em separado:
remove temporariamente $e$, realiza uma BFS ou DFS a partir de $u$
e verifica se $v$ continua alcançável. Caso contrário, $e$ é ponte.
Depois, a aresta é restaurada. Uma busca tem custo $\ordem{n+m}$;
como isso pode ocorrer para cada uma das $m$ arestas, o custo total é
\begin{equation}
T_{\text{naive}}(n,m)=\ordem{m(n+m)}.
\end{equation}
Uma otimização útil interrompe a busca assim que $v$ é alcançado
(\emph{early exit}). A resposta não muda, mas o tempo médio pode cair.

\subsection{Identificação de pontes por Tarjan}
O algoritmo de Tarjan executa uma única DFS e mantém os seguintes valores:
$disc[u]$, a ordem de descoberta de $u$, e $low[u]$, o menor tempo de
descoberta de um ancestral alcançável a partir da subárvore de $u$
por arestas de árvore e, no máximo, uma aresta de retorno.
Para uma aresta de árvore $(u,v)$, com $u$ pai de $v$, vale:
\begin{equation}\label{eq:tarjan}
\boxed{(u,v)\text{ é ponte}\quad\Longleftrightarrow\quad low[v]>disc[u].}
\end{equation}
Se $low[v]\le disc[u]$, há um caminho alternativo que fecha um ciclo;
se $low[v]>disc[u]$, a subárvore de $v$ depende daquela aresta
para se conectar aos ancestrais. Como cada vértice e aresta é visitado
um número constante de vezes, a complexidade é $\ordem{n+m}$.

\subsection{Algoritmo de Fleury}
Fleury constrói um percurso consumindo uma aresta a cada passo. A regra
central é evitar atravessar uma ponte do \emph{grafo residual} sempre que
houver outra aresta incidente disponível. Se o grau ativo do vértice
corrente é um, essa aresta é obrigatória.

\begin{lstlisting}[caption={Pseudocódigo do algoritmo de Fleury.}]
FLEURY(G, estrategia):
    se G nao satisfaz as condicoes de Euler:
        retornar NAO_EXISTE
    u = vertice impar, se houver; senao vertice com arestas
    caminho = [u]
    enquanto houver arestas:
        se grauAtivo(u) == 1:
            e = unica aresta incidente
        senao:
            e = primeira aresta incidente que nao e ponte
                segundo a estrategia escolhida
        remover e do grafo residual
        u = outra extremidade de e
        adicionar u a caminho
    retornar caminho
\end{lstlisting}

Uma execução de Fleury realiza $m$ passos. Se Tarjan é recalculado
integralmente a cada decisão, o limite superior é $\ordem{m(n+m)}$.
Se o método \naive{} identifica \emph{todas} as pontes em cada passo,
o limite chega a $\ordem{m^2(n+m)}$. Entretanto, ao testar apenas
as arestas candidatas de Fleury, normalmente basta uma ou duas buscas
por passo. Nesse caso, o limite superior volta a ser $\ordem{m(n+m)}$.

\begin{table}[H]
\centering\small
\caption{Comparação das complexidades assintóticas.}
\label{tab:complexidade}
\begin{tabular}{@{}p{7.0cm}p{5.1cm}@{}}
\toprule
Procedimento & Limite superior \\
\midrule
Busca BFS/DFS & $\ordem{n+m}$ \\
Todas as pontes: método naïve & $\ordem{m(n+m)}$ \\
Todas as pontes: Tarjan & $\ordem{n+m}$ \\
Fleury + Tarjan recalculado & $\ordem{m(n+m)}$ \\
Fleury + naïve sob demanda & $\ordem{m(n+m)}$ \\
Fleury + naïve para todas as arestas & $\ordem{m^2(n+m)}$ \\
\bottomrule
\end{tabular}
\end{table}

\section{Implementação em Java}
O projeto foi organizado no pacote \texttt{grafos}, sem bibliotecas externas.
A Tabela~\ref{tab:arquivos} apresenta os arquivos centrais da implementação.

\begin{table}[H]
\centering\small
\caption{Principais classes do projeto.}
\label{tab:arquivos}
\begin{tabular}{@{}p{5cm}p{7.2cm}@{}}
\toprule
Classe & Responsabilidade \\
\midrule
\texttt{Graph} & Estrutura do grafo e remoção/restauração de arestas. \\
\texttt{Reach} & BFS para testes de alcançabilidade. \\
\texttt{NaiveBridgeFinder} & Identificação naïve de pontes. \\
\texttt{TarjanBridgeFinder} & DFS iterativa com vetores de descoberta e retorno. \\
\texttt{EulerChecker} & Verificação das condições eulerianas. \\
\texttt{BridgeStrategy} e \texttt{Strategies} & Integração dos detectores ao Fleury. \\
\texttt{Fleury} & Construção de caminhos ou circuitos eulerianos. \\
\texttt{GraphGenerator} & Geração determinística de grafos aleatórios. \\
\texttt{EulerPathValidator} & Validação independente das arestas percorridas. \\
\texttt{Benchmark}, \texttt{SelfTest}, \texttt{StressTest} & Experimentos e testes automatizados. \\
\bottomrule
\end{tabular}
\end{table}

\subsection{Representação do grafo}
O grafo usa listas de adjacência em vetores contíguos. Cada aresta
possui identificador estável; os graus ativos e as posições das arestas
são atualizados quando ocorre uma remoção lógica. A restauração desfaz
as alterações da última remoção, permitindo repetir testes de pontes
sem modificar permanentemente o grafo. A DFS de Tarjan é iterativa,
com pilha explícita, para evitar estouro da pilha de chamadas com
100.000 vértices.

\subsection{Pseudocódigo dos detectores}
\begin{lstlisting}[caption={Detecção de todas as pontes por força bruta.}]
PONTES_NAIVE(G):
    pontes = conjunto vazio
    para cada aresta e=(u,v):
        remover_temporariamente(e)
        conectados = BFS_ALCANCA(G,u,v)
        restaurar(e)
        se nao conectados: adicionar e a pontes
    retornar pontes
\end{lstlisting}

\begin{lstlisting}[caption={Regra de Tarjan em uma DFS; no Java a busca é iterativa.}]
TARJAN(u):
    disc[u] = low[u] = tempo++
    para cada aresta e=(u,v) diferente da aresta do pai:
        se v nao foi visitado:
            pai[v] = u
            TARJAN(v)
            low[u] = min(low[u], low[v])
            se low[v] > disc[u]: marcar e como ponte
        senao:
            low[u] = min(low[u], disc[v])
\end{lstlisting}

\subsection{Estratégias utilizadas pelo Fleury}
Foram implementadas quatro estratégias para interpretar tanto a comparação
principal quanto o custo de computar pontes sob demanda:
\begin{itemize}
\item \texttt{tarjan}: recalcula as pontes do grafo residual por DFS;
\item \texttt{naive}: percorre toda a componente para testar cada candidata;
\item \texttt{naive-early}: interrompe a BFS ao encontrar o destino;
\item \texttt{naive-all}: identifica todas as pontes por força bruta em cada decisão.
\end{itemize}
As duas abordagens requeridas são Tarjan e naïve; as variantes adicionais
servem somente para explicar diferenças de desempenho.

\subsection{Exemplo de funcionamento}
Considere um triângulo com vértices $0,1,2$ e o caminho $2-3-4$.
As arestas $\{2,3\}$ e $\{3,4\}$ são pontes porque não participam de
ciclos. Os únicos vértices de grau ímpar são $2$ e $4$; portanto, há
um caminho euleriano aberto.

\begin{figure}[H]
\centering
\begin{tikzpicture}[scale=1.05, every node/.style={font=\sffamily\small}]
  \coordinate (a) at (0,0); \coordinate (b) at (2,0);
  \coordinate (c) at (1,1.7); \coordinate (d) at (3.6,1.7);
  \coordinate (e) at (5.6,1.7);
  \draw[thick] (a)--(b)--(c)--cycle;
  \draw[very thick,dashed] (c)--(d)--(e);
  \foreach \p/\lab in {a/0,b/1,c/2,d/3,e/4}{
    \filldraw[fill=white,draw=black,thick] (\p) circle (0.19);
    \node at (\p) {\lab};
  }
  \node[above] at (2.3,2.05) {pontes};
\end{tikzpicture}
\caption{Exemplo de grafo semieuleriano: as arestas tracejadas são pontes.}
\label{fig:exemplo}
\end{figure}

Um percurso correto de Fleury é
$2\rightarrow0\rightarrow1\rightarrow2\rightarrow3\rightarrow4$.
As cinco arestas aparecem exatamente uma vez. Fleury evita atravessar
$2-3$ enquanto as arestas do triângulo ainda precisam ser percorridas.


\noindent\textbf{Grafo da Figura 1 do enunciado.}
O programa também inclui o grafo original de exemplo com nove vértices
($A$ a $I$) e quinze arestas. A rotina \texttt{Main demo} encontra uma
única ponte, $\{F,G\}$, tanto pelo método \naive{} quanto por Tarjan.
Os vértices de grau ímpar são $A$ e $G$, e Fleury encontra o caminho:
\[
A\to D\to F\to C\to B\to A\to E\to B\to D\to C
\to E\to F\to G\to I\to H\to G.
\]
O validador independente confirma que todas as quinze arestas foram
percorridas exatamente uma vez, sem alterar o grafo de entrada.

\subsection{Geração e validação dos grafos}
O gerador começa com um ciclo hamiltoniano aleatorizado e acrescenta
ciclos curtos, garantindo graus pares e conexidade. A remoção de uma
aresta de um grafo euleriano produz um grafo semieuleriano; a remoção
de duas arestas sem extremos comuns produz quatro vértices ímpares,
caracterizando um grafo não euleriano. As sementes são fixas para permitir
reprodução. Há ainda um gerador de grafos conexos com árvore aleatória e
arestas adicionais, útil para testar instâncias com muitas pontes.

A validação confere a quantidade de vértices do caminho, a adjacência
dos consecutivos e o uso de cada aresta exatamente uma vez. Os testes
principais registrados realizaram 31.385 verificações sem falhas; uma
bateria adicional de estresse verificou grafos com até 100.000 vértices.

\section{Metodologia experimental}
\subsection{Ambiente e procedência dos dados}
Os resultados preservados nos CSVs do projeto consolidado foram obtidos
em ambiente Linux Ubuntu 24.04, com processador Intel Xeon de 2,30 GHz,
7,8 GiB de RAM e OpenJDK 21. A medição empregou
\texttt{System.nanoTime()}, aquecimento da JVM e verificação posterior
dos caminhos. Os tempos são apresentados com as unidades originais;
não correspondem a um novo benchmark executado na máquina de entrega.

\subsection{Planejamento dos experimentos}
Foram utilizados quatro tamanhos: $n\in\{100,1000,10000,100000\}$.
No experimento principal, o grau médio foi aproximadamente quatro
($m\approx2n$). O número de repetições diminuiu com o tamanho,
por causa do crescimento do tempo de Fleury:
\begin{itemize}
\item E1: Fleury com Tarjan, naïve e naïve com parada antecipada;
\item E2: Fleury com naïve recalculando todas as pontes;
\item E3: comparação isolada dos algoritmos de detecção de pontes;
\item E4: sensibilidade de Fleury ao aumento de densidade.
\end{itemize}

Os números representam médias aritméticas das execuções concluídas.
Os resultados de 100.000 vértices incluem apenas duas repetições para
grafos eulerianos e uma para cada uma das outras classes, o que limita
a precisão estatística. Os dados originais permitem repetir e ampliar
os testes em outras máquinas.

\section{Resultados}
\subsection{E1: Fleury em grafos eulerianos e semieulerianos}
A Tabela~\ref{tab:principal} reúne os tempos médios de Fleury nas duas
classes para as quais existe caminho euleriano. A versão \texttt{naive-early}
teve o menor tempo nos grafos desta família.

\begin{table}[H]
\centering\small
\caption{Fleury: tempos médios dos experimentos históricos.}
\label{tab:principal}
\resizebox{\textwidth}{!}{%
\begin{tabular}{@{}llrrrrr@{}}
\toprule
Classe & $n$ & $m$ aprox. & Reps. & Tarjan & Naïve & Naïve (early exit) \\
\midrule
Euleriano & 100 & 202 & 200 & 0,20 ms & 0,17 ms & 0,069 ms \\
 & 1.000 & 2.003 & 50 & 20,5 ms & 12,8 ms & 4,51 ms \\
 & 10.000 & 20.002 & 10 & 3,36 s & 1,44 s & 0,487 s \\
 & 100.000 & 200.004 & 2 & 18,7 min & 3,96 min & 75,5 s \\
\midrule
Semieuleriano & 100 & 202 & 200 & 0,21 ms & 0,15 ms & 0,068 ms \\
 & 1.000 & 2.002 & 50 & 20,5 ms & 12,7 ms & 4,52 ms \\
 & 10.000 & 20.001 & 10 & 3,26 s & 1,45 s & 0,486 s \\
 & 100.000 & 200.002 & 1 & 18,4 min & 3,95 min & 77,4 s \\
\bottomrule
\end{tabular}}
\end{table}

\begin{figure}[H]
\centering
\begin{tikzpicture}
\begin{loglogaxis}[
 width=0.87\textwidth,height=6.1cm,
 xlabel={Número de vértices $n$},ylabel={Tempo médio (ms)},
 xtick={100,1000,10000,100000},
 xticklabels={$10^2$,$10^3$,$10^4$,$10^5$},
 grid=both,legend pos=north west,
 legend style={font=\footnotesize},tick label style={font=\footnotesize}]
\addplot+[black,mark=square*,thick] coordinates
 {(100,0.19932) (1000,20.53424) (10000,3357.45606) (100000,1123317.80107)};
\addlegendentry{Fleury + Tarjan}
\addplot+[black!65,mark=*,thick,dashed] coordinates
 {(100,0.16771) (1000,12.78916) (10000,1438.20416) (100000,237500.40057)};
\addlegendentry{Fleury + Naïve}
\addplot+[black!45,mark=triangle*,thick,dotted] coordinates
 {(100,0.06948) (1000,4.50621) (10000,486.61040) (100000,75496.74912)};
\addlegendentry{Fleury + Naïve early exit}
\end{loglogaxis}
\end{tikzpicture}
\caption{Crescimento do tempo de Fleury em grafos eulerianos (eixos logarítmicos).}
\label{fig:fleury}
\end{figure}

Nos grafos não eulerianos, o teste inicial de paridade encontrou quatro
vértices ímpares e o percurso foi corretamente rejeitado, sem necessidade
de executar a escolha de arestas de Fleury. A Tabela~\ref{tab:naoeuleriano}
registra os tempos dessa checagem, que não representam travessias completas.

\begin{table}[H]
\centering\small
\caption{Grafos não eulerianos: tempo médio de rejeição do percurso ($\mu$s).}
\label{tab:naoeuleriano}
\begin{tabular}{@{}rrrrr@{}}
\toprule
$n$ & Repetições & Tarjan & Naïve & Naïve (early exit) \\
\midrule
100 & 200 & 1,05 & 0,77 & 1,53 \\
1.000 & 50 & 1,10 & 0,82 & 0,76 \\
10.000 & 10 & 19,48 & 16,53 & 17,42 \\
100.000 & 1 & 249,10 & 353,44 & 333,39 \\
\bottomrule
\end{tabular}
\end{table}

Esses tempos são pequenos porque o programa interrompe a execução
após a verificação dos graus. Para os menores grafos, valores inferiores
a poucos microssegundos devem ser interpretados com cautela devido à
resolução do relógio e às oscilações da JVM.

\subsection{E2: Fleury com força bruta sobre todas as pontes}
A variante \texttt{naive-all} calcula novamente o conjunto completo de
pontes em cada passo de decisão. Esse procedimento é substancialmente
mais caro que testar apenas as candidatas.

\begin{table}[H]
\centering\small
\caption{Fleury com naïve aplicado a todas as arestas do grafo residual.}
\label{tab:naiveall}
\begin{tabular}{@{}llrr@{}}
\toprule
Classe & $n$ & Naïve-todas & Tarjan \\
\midrule
Euleriano & 100 & 23,3 ms & 0,20 ms \\
Euleriano & 1.000 & 21,5 s & 20,9 ms \\
Euleriano & 10.000 & $>120$ s (limite) & 3,26 s \\
Semieuleriano & 100 & 17,4 ms & 0,19 ms \\
Semieuleriano & 1.000 & 21,6 s & 20,8 ms \\
Semieuleriano & 10.000 & $>120$ s (limite) & 3,18 s \\
\bottomrule
\end{tabular}
\end{table}

O limite de 120 segundos foi atingido nas instâncias de 10.000 vértices
antes da conclusão do percurso. Portanto, os valores registrados como
limite de tempo não são médias completas nem devem ser interpretados
como tempos exatos do algoritmo.

\subsection{E3: detecção isolada de todas as pontes}
Para comparar os detectores sem a interferência de Fleury, foram usados
grafos conexos aleatórios com grau médio três, contendo pontes.

\begin{table}[H]
\centering\small
\caption{Detecção completa de pontes: tempos médios históricos.}
\label{tab:pontes}
\resizebox{\textwidth}{!}{%
\begin{tabular}{@{}rrrrrrr@{}}
\toprule
$n$ & $m$ & Reps. & Pontes & Tarjan & Naïve & Naïve (early exit) \\
\midrule
100 & 150 & 100 & 20 & 31,1 $\mu$s & 0,40 ms & 0,28 ms \\
1.000 & 1.500 & 50 & 212 & 0,24 ms & 32,6 ms & 12,9 ms \\
10.000 & 15.000 & 10 & 2.123 & 0,64 ms & 3,00 s & 1,26 s \\
100.000 & 150.000 & 1 & 21.122 & 16,7 ms & 7,59 min & 3,25 min \\
\bottomrule
\end{tabular}}
\end{table}

\begin{figure}[H]
\centering
\begin{tikzpicture}
\begin{loglogaxis}[
 width=0.86\textwidth,height=6.2cm,
 xlabel={Número de vértices $n$},ylabel={Tempo médio (ms)},
 xtick={100,1000,10000,100000},
 xticklabels={$10^2$,$10^3$,$10^4$,$10^5$},
 grid=both,legend pos=north west,
 legend style={font=\footnotesize},tick label style={font=\footnotesize}]
\addplot+[black,mark=square*,thick] coordinates
 {(100,0.03113) (1000,0.23541) (10000,0.64426) (100000,16.73522)};
\addlegendentry{Tarjan}
\addplot+[black!65,mark=*,thick,dashed] coordinates
 {(100,0.40175) (1000,32.64271) (10000,2999.69903) (100000,455275.81582)};
\addlegendentry{Naïve}
\addplot+[black!45,mark=triangle*,thick,dotted] coordinates
 {(100,0.27789) (1000,12.89392) (10000,1264.72704) (100000,194708.68781)};
\addlegendentry{Naïve early exit}
\end{loglogaxis}
\end{tikzpicture}
\caption{Detecção isolada de pontes, em escala logarítmica.}
\label{fig:pontes}
\end{figure}

Para 100.000 vértices, Tarjan levou 16,7 ms, enquanto o método
\naive{} levou aproximadamente 7,59 minutos. A razão obtida no
registro histórico foi de aproximadamente 27 mil vezes. Isso é
compatível com os limites $\ordem{n+m}$ e $\ordem{m(n+m)}$ em grafos
esparsos, embora uma única execução não permita estimar a variabilidade.

\subsection{E4: efeito da densidade}
Para $n=1.000$ vértices eulerianos, aumentou-se o grau médio de
aproximadamente quatro para 32, mantendo a classificação euleriana.

\begin{table}[H]
\centering\small
\caption{Fleury em grafos eulerianos de 1.000 vértices, variando o grau médio.}
\label{tab:densidade}
\begin{tabular}{@{}rrrrr@{}}
\toprule
Grau médio & $m$ & Tarjan & Naïve & Naïve (early exit) \\
\midrule
4 & 2.003 & 28,4 ms & 17,6 ms & 6,22 ms \\
8 & 4.002 & 105 ms & 55,4 ms & 16,8 ms \\
16 & 8.002 & 359 ms & 157 ms & 37,4 ms \\
32 & 16.003 & 1,17 s & 412 ms & 66,6 ms \\
\bottomrule
\end{tabular}
\end{table}

\begin{figure}[H]
\centering
\begin{tikzpicture}
\begin{loglogaxis}[
 width=0.86\textwidth,height=6.1cm,
 xlabel={Número de arestas $m$},ylabel={Tempo médio (ms)},
 grid=both,legend pos=north west,
 legend style={font=\footnotesize},tick label style={font=\footnotesize}]
\addplot+[black,mark=square*,thick] coordinates
 {(2003,28.44667) (4002,105.37502) (8002,359.02060) (16003,1172.96196)};
\addlegendentry{Tarjan}
\addplot+[black!65,mark=*,thick,dashed] coordinates
 {(2003,17.63670) (4002,55.42889) (8002,157.33488) (16003,412.18028)};
\addlegendentry{Naïve}
\addplot+[black!45,mark=triangle*,thick,dotted] coordinates
 {(2003,6.21997) (4002,16.79876) (8002,37.44380) (16003,66.63545)};
\addlegendentry{Naïve early exit}
\end{loglogaxis}
\end{tikzpicture}
\caption{Impacto da densidade no tempo de Fleury, mantendo $n=1.000$.}
\label{fig:densidade}
\end{figure}

À medida que a densidade aumentou, o percentual de passos forçados
caiu de aproximadamente 50\% para 6,2\%. Mesmo assim, a verificação
local das arestas candidatas permaneceu mais rápida que o recálculo
completo de todas as pontes por Tarjan neste conjunto de grafos.

\section{Discussão}
\subsection{Por que Tarjan ganha na detecção completa?}
Na detecção de todas as pontes, Tarjan usa uma única DFS com marcações
de descoberta e valores \emph{low}. O método \naive{} repete uma busca
para cada aresta, causando crescimento muito mais rápido quando o
número de vértices e arestas aumenta. O experimento E3 confirma
qualitativamente essa diferença.

\subsection{Por que o naïve local pode vencer dentro de Fleury?}
Durante Fleury, somente as arestas incidentes ao vértice corrente
precisam ser analisadas. Se a primeira aresta não é ponte, uma única
BFS decide o passo. A estratégia Tarjan, por sua vez, recalcula o
conjunto completo de pontes do grafo residual. Nos grafos usados,
cerca de metade dos passos de Fleury já era forçada (grau ativo um),
e nos demais a primeira candidata quase sempre era adequada.
Isso explica por que o \naive{} sob demanda foi mais rápido nessas
instâncias, sem contradizer a superioridade assintótica de Tarjan
para identificação completa de pontes.

\subsection{Limitações da comparação}
Os resultados devem ser interpretados sob quatro restrições:
\begin{itemize}
\item os geradores produzem famílias de grafos específicas, não uma amostra uniforme de todos os grafos simples;
\item as medições de 100.000 vértices possuem poucas repetições;
\item os tempos dependem de hardware, JVM, aquecimento, coleta de lixo e ordenação das arestas;
\item as estratégias realizam quantidades diferentes de trabalho em cada consulta, embora todas obedeçam à regra de Fleury.
\end{itemize}
O experimento de densidade amplia a análise, mas não elimina essas
restrições. Um trabalho futuro poderia comparar outras topologias,
mais sementes, diferentes máquinas e o algoritmo de Hierholzer,
que encontra caminhos eulerianos em tempo linear no número de arestas.

\section{Conclusão}
Foram implementados dois métodos de detecção de pontes e quatro
estratégias de integração ao algoritmo de Fleury. A aplicação inclui
geradores aleatórios, verificador independente de caminhos eulerianos
e testes automatizados, com casos pequenos e instâncias de até
100.000 vértices.

Na identificação isolada de todas as pontes, Tarjan apresentou o melhor
desempenho, em conformidade com sua complexidade linear. Dentro de
Fleury, entretanto, os testes locais do método \naive{} foram mais
rápidos para as famílias esparsas e eulerianas avaliadas, porque evitam
calcular pontes desnecessárias. A variante que recalcula todas as pontes
por força bruta a cada passo mostrou-se inviável em grafos grandes.

Portanto, a escolha do método de identificação de pontes deve levar
em conta \emph{como e quantas vezes} ele será consultado. Tarjan é
recomendável para a identificação global de pontes; o teste local pode
ser vantajoso em aplicações que precisam decidir sobre poucas arestas
por vez, como na implementação de Fleury analisada.

\section*{Referências}
\addcontentsline{toc}{section}{Referências}
\begin{enumerate}[label={[\arabic*]},leftmargin=1.15cm]
\item TARJAN, R. E. A note on finding the bridges of a graph.\
\emph{Information Processing Letters}, v. 2, n. 6, p. 160--161, 1974.\
\url{https://doi.org/10.1016/0020-0190(74)90003-9}.
\item FLEURY, M. Deux problèmes de géométrie de situation.\
\emph{Journal de mathématiques élémentaires}, 2e série, v. 2, p. 257--261, 1883.
\item HIERHOLZER, C. Über die Möglichkeit, einen Linienzug ohne Wiederholung und ohne Unterbrechung zu umfahren.\
\emph{Mathematische Annalen}, v. 6, p. 30--32, 1873.
\item CORMEN, T. H. et al. \emph{Introduction to Algorithms}. 3. ed. MIT Press, 2009.
\item WEST, D. B. \emph{Introduction to Graph Theory}. 2. ed. Prentice Hall, 2001.
\end{enumerate}

\end{document}
